% Generated draft: requires mathematical and source review.
\documentclass[11pt]{article}
\usepackage[margin=25mm]{geometry}
\usepackage{fontspec}
\setmainfont{Latin Modern Roman}
\usepackage{amsmath,amssymb,mathtools,unicode-math}
\setmathfont{Latin Modern Math}
\usepackage{enumitem,needspace,xurl,hyperref}
\setlength{\parindent}{0pt}
\setlength{\parskip}{5pt}
\setlength{\emergencystretch}{3em}
\newcommand{\N}{\mathbb{N}}
\newcommand{\Z}{\mathbb{Z}}
\newcommand{\Q}{\mathbb{Q}}
\newcommand{\R}{\mathbb{R}}
\newcommand{\C}{\mathbb{C}}
\newcommand{\F}{\mathbb{F}}
\newcommand{\A}{\mathbb{A}}
\newcommand{\PP}{\mathbb{P}}
\newcommand{\E}{\mathbb{E}}
\newcommand{\eps}{\varepsilon}
\newcommand{\dd}{\,\mathrm{d}}
\begin{document}
\section*{UnsolvedMath ID: 701753}
\textbf{Source problem code:} catalogue entry
\subsection{Definitions and mathematical statement}
A group \(G\) is periodic if every element has finite order. An involution is an element of order \(2\). The condition on subgroups of order \(4\) is
\[
\forall V\leq G,\qquad |V|=4\ \Longrightarrow\ V\text{ is cyclic}.
\]
A finite group \(F\) is a finite Frobenius group if it has a nontrivial proper subgroup \(C\) such that
\[
C\cap x^{-1}Cx=\{1\}\qquad\text{for every }x\in F\setminus C.
\]
Using the usual subgroup-containment interpretation, \(G\) is saturated with finite Frobenius groups if every finite subgroup \(H\leq G\) is contained in a finite subgroup \(F\leq G\) that is a Frobenius group. The supplied record does not explicitly define saturation.

The source asks: if \(G\) is a periodic group saturated with finite Frobenius groups, contains an involution, and has no noncyclic subgroup of order \(4\), must \(G\) be a Frobenius group in the sense of Kourovka Notebook Problem 6.53?

The referenced definition in Problem 6.53 is not supplied. Under the convention that an arbitrary group is Frobenius when it contains a nontrivial proper malnormal subgroup, the precise conclusion would be
\[
\exists C\leq G:\quad \{1\}\subsetneq C\subsetneq G
\quad\text{and}\quad
\forall g\in G\setminus C,\quad C\cap g^{-1}Cg=\{1\}.
\]
Thus the conditional formulation asks whether every group satisfying the stated hypotheses has such a subgroup \(C\). This convention must be checked against Problem 6.53 before treating it as the source's exact intended conclusion. No finiteness or local finiteness assumption on \(G\) is added.
\subsection{Short English statement}
Must every periodic group that is saturated with finite Frobenius groups, contains an involution, and has only cyclic subgroups of order \(4\) be a Frobenius group? The source refers to Problem 6.53 for the intended meaning of Frobenius group; that definition is absent, so the explicit malnormal-subgroup formulation remains conditional.
\subsection{Sources}
\begin{itemize}
\item Supplied local UnsolvedMath registry record, permanent ID \(701753\), exact source code \texttt{catalogue entry}, titled \emph{Kourovka Notebook Problem 20.95}. The record labels the problem \emph{open}; this is supplied registry status, not an independently verified mathematical assessment.
\item Attributed to A. I. Sozutov, proposed in 2022. The supplied source identifies \emph{The Kourovka Notebook (Unsolved Problems in Group Theory)}, 20th issue (2022), Problem 20.95, and the 21st issue edition (2026), \texttt{arXiv:1401.0300v48}. The collection description names E. I. Khukhro and V. D. Mazurov as editors.
\item The statement refers to Kourovka Notebook Problem 6.53 for the meaning of Frobenius group. Its text is not included in the supplied evidence, and the cited editions have not been independently consulted.
\end{itemize}
\end{document}
